\begin{mistake}{2026-05-25}{04}

\begin{problemcontent}
计算不定积分：
\( I = \int \frac{x^4 + 1}{1 + x^6} \mathrm{d}x \)
\end{problemcontent}

\begin{solutioncontent}
对分母利用立方和公式进行因式分解：
\[
1 + x^6 = 1^3 + (x^2)^3 = (1 + x^2)(1 - x^2 + x^4)
\]
观察到分母出现的长因式 \(x^4 - x^2 + 1\) 与分子 \(x^4 + 1\) 仅相差一项 \(-x^2\)，因此对分子通过加减项进行拼凑拆分：
\[\begin{aligned}
\frac{x^4 + 1}{1 + x^6} &= \frac{(x^4 - x^2 + 1) + x^2}{(1 + x^2)(x^4 - x^2 + 1)} \\
&= \frac{x^4 - x^2 + 1}{(1 + x^2)(x^4 - x^2 + 1)} + \frac{x^2}{1 + (x^3)^2} \\
&= \frac{1}{1 + x^2} + \frac{x^2}{1 + (x^3)^2}
\end{aligned}\]
分别对两项进行积分，第二项使用凑微分法 \(x^2 \mathrm{d}x = \frac{1}{3}\mathrm{d}(x^3)\)：
\[\begin{aligned}
I &= \int \frac{1}{1 + x^2} \mathrm{d}x + \int \frac{x^2}{1 + (x^3)^2} \mathrm{d}x \\
&= \arctan x + \frac{1}{3} \int \frac{1}{1 + (x^3)^2} \mathrm{d}(x^3) \\
&= \arctan x + \frac{1}{3} \arctan(x^3) + C
\end{aligned}\]
\end{solutioncontent}

\mistakeknowledge{
\begin{itemize}
  \item \textbf{核心公式}：立方和公式 \(a^3+b^3=(a+b)(a^2-ab+b^2)\)。
  \item \textbf{解题技巧}：面对高次有理分式 \(x^n \pm 1\)，首选通过因式分解提取公因式，并与分子进行凑项比对，避免复杂的待定系数法部分分式展开。
  \item \textbf{关联知识}：基础积分公式 \(\int \frac{1}{1+u^2}\mathrm{d}u = \arctan u + C\)。
\end{itemize}}

\end{mistake}
